% intro.tex (v10)
The active layer is the ground above permafrost that thaws in summer and refreezes in winter. Its thickness (ALT) controls the depth of seasonal thaw, the stability of infrastructure and the amount of soil carbon exposed to decomposition\cite{Hjort2018,Schuur2015,Biskaborn2019}. ALT changes from year to year with the summer heat input, and the year-to-year signal is part of what monitoring programmes such as the Circumpolar Active Layer Monitoring network record\cite{Nelson1997,Streletskiy2025}. Observations, however, are sparse and uneven: most measurements come from a few dense campaigns and long-term grids, and many sites are visited only once\cite{Moore2025,Veremeeva2025}.

Several 1 km ALT products cover the northern permafrost region. Process-based estimates such as ESA CCI Permafrost run a heat-conduction model with remotely sensed land-surface temperature\cite{Westermann2024,Westermann2025}. Statistical and machine-learning (ML) products relate site observations to climate, terrain and soil covariates\cite{Aalto2018,Ran2022a,Liu2024,Wei2026}. These products are typically trained on multi-year mean conditions or on site observations from many years pooled together, and they are evaluated with different site sets and validation schemes. A product value therefore cannot be compared directly with an observation of a given cell and year, and products differ by tens of centimetres in the same cell. Validation design adds a further uncertainty: random cross-validation of spatially clustered observations overstates the accuracy of ML models for new locations\cite{Roberts2017,Ploton2020,MeyerPebesma2022,Wadoux2021}.

Physical models offer a complementary route. The Stefan solution, $\mathrm{ALT}=E\sqrt{\mathrm{TDD}}$, links thaw depth to the thawing degree-days (TDD) of the season through a single coefficient $E$ that depends on soil thermal properties and moisture\cite{Stefan1891,Riseborough2008}. It reproduces the seasonal course of thaw and responds to the climate of each year, but one coefficient per region leaves systematic residuals that vary with soil organic carbon and terrain. Combining physical models with ML has improved environmental predictions when the physical model carries the dominant response and the ML model learns what the physics misses\cite{Willard2022}. For ALT, process-based and ML estimates have been combined for Alaska\cite{Gautam2025}, but it is not known how such a combination behaves when the input climate is matched to the observation year, when it is transferred to a region without labels, or how many labels a new region needs.

Here we address these questions with a cell-year design: every ALT observation is a 1 km cell in a given year, and every climate input is computed for that year from ERA5-Land, accumulated to the observation day\cite{munozsabater2021_era5_land}. Using \nncyAll{} cell-year observations from 2010--2024 in Alaska, the Lena Delta and north-western Canada, we compare the Stefan model, direct ML and a physics-residual workflow in which gradient boosting learns the residual of a recalibrated Stefan model. The models are evaluated under five validation designs, along label-number ladders within and across regions, in external permafrost-monitoring regions and against six published products in the same cells and years. We report (1) the effect of the year-specific climate input, (2) the gain of the physics-residual model over the Stefan model and its conditions, (3) the number of labels needed in a new region and (4) the differences from published products, and we produce year-labelled 1 km maps for 2011--2024.
